%% LyX 2.4.4 created this file.  For more info, see https://www.lyx.org/.
%% Do not edit unless you really know what you are doing.
\documentclass[english]{article}
\usepackage[T1]{fontenc}
\usepackage[latin9]{inputenc}
\usepackage{units}
\usepackage{enumitem}
\usepackage{amsmath}
\usepackage{amsthm}
\usepackage{amssymb}
\usepackage{cancel}
\usepackage{geometry}
\geometry{verbose,tmargin=1in,bmargin=1in,lmargin=1in,rmargin=1in}

\makeatletter
%%%%%%%%%%%%%%%%%%%%%%%%%%%%%% Textclass specific LaTeX commands.
\numberwithin{equation}{section}
\numberwithin{figure}{section}
\newlength{\lyxlabelwidth}      % auxiliary length 

%%%%%%%%%%%%%%%%%%%%%%%%%%%%%% User specified LaTeX commands.
\usepackage{fancyhdr}
\usepackage{lastpage}
\pagestyle{fancy}
\usepackage{enumitem}
\usepackage{ifthen}
%sets math fonts to be more readable
\everymath=\expandafter{\the\everymath\displaystyle}
%\DeclareMathSizes{10}{10}{10}{10}
\usepackage{multicol}

\usepackage{tikz}
\usetikzlibrary{plotmarks}
\usepackage{pgfplots}
\pgfplotsset{compat=newest}

\usepackage{calc}
\usepackage[nomessages]{fp}

\usepackage{advdate}    % Advancing/saving dates
\usepackage{datetime}   % Dates formatting
\usepackage{datenumber} % Counters for dates
% Added by lyx2lyx
\setlength{\parskip}{\smallskipamount}
\setlength{\parindent}{0pt}

\makeatother

\usepackage{babel}
\begin{document}
\renewcommand{\author}{Dr.\,James Ashe}

\newcommand{\classNum}{131}

\newcommand{\classSec}{B}

\newcommand{\nameType}{}

\newcommand{\examType}{Exam 2 Review}

\renewcommand{\date}{\formatdate{30}{9}{2013}}

%%First page's header%%%%%%%%%%%%%%%%%%%%%%
\fancypagestyle{firststyle} %%header settings for first pagr
{
	\fancyhead[L]{MTH \classNum{}(\classSec{}) \author{}\\
	\textbf{\examType{}} \date{}}
	\fancyhead[C]{\textbf{\nameType}}
	\setlength{\headsep}{14pt}
}
%%%%%%%%%%%%%%%%%%%%%%%%%%%%%%%%%%%%%%%%%%

%%Next page's header%%%%%%%%%%%%%%%%%%%%%%%
\setlist{nolistsep}
\lhead{\classNum{}(\classSec{})} 
\chead{\textbf{\nameType}} 
\cfoot{\thepage{}~of~\pageref{LastPage}} 
\setlength{\headsep}{5pt} 
%%%%%%%%%%%%%%%%%%%%%%%%%%%%%%%%%%%%%%%%%%%

%%Various settings; see the preamble for other settings%%%%%
%%\setenumerate[0]{leftmargin=12pt,itemindent=0pt,labelwidth=0pt}
\renewcommand{\labelenumi}{\textbf{\arabic{enumi}.}}
\renewcommand{\labelenumii}{\textbf{\alph{enumii}.}}
\thispagestyle{firststyle}
%%%%%%%%%%%%%%%%%%%%%%%%%%%%%%%%%%%%%%%%%%
\newcommand{\xmax}{10}
\newcommand{\xmin}{-10}
\newcommand{\xstep}{0} %must be xmin+n
\newcommand{\ymax}{10}
\newcommand{\ymin}{-10}
\newcommand{\ystep}{0} %must be ymin+n

\newcommand{\Dspace}{\vspace{1cm}}

\small\textbf{Solve a radical. ex. }$\sqrt{x-1}+2=5$\textbf{.}
\begin{enumerate}
\item Get the radical by itself, 
\begin{eqnarray*}
\sqrt{x-1}+2 & = & 5\\
\sqrt{x-1} & = & 3
\end{eqnarray*}
\item Use a square root, simplify, and solve. 
\begin{eqnarray*}
\left(\sqrt{x-1}\right)^{2} & = & 3^{2}\\
x-1 & = & 9\\
x & = & 10
\end{eqnarray*}
\item Check! $\sqrt{10-1}=9\mbox{\ensuremath{\checkmark}}$
\end{enumerate}
\textbf{Solve an equation with rational exponents. ex. }$(x-7)^{\nicefrac{3}{2}}-64=0$\textbf{.}
\begin{enumerate}
\item Get the exponent by itself.
\begin{eqnarray*}
(x-7)^{\nicefrac{3}{2}} & = & 64.
\end{eqnarray*}
\item Get rid of the exponent, using the rule $(b^{n})^{m}=b^{nm}$. Raise
both sides by the reciprocal of the exponent (``flip it''). 
\begin{eqnarray*}
\left((x-7)^{\nicefrac{3}{2}}\right)^{\nicefrac{2}{3}} & = & (64)^{\nicefrac{2}{3}}.\\
x-7 & = & 64^{\nicefrac{2}{3}}\\
x & = & 16+7\\
x & = & 23
\end{eqnarray*}
\item Check. $(23-7)^{\nicefrac{3}{2}}=64.\checkmark$
\end{enumerate}
\textbf{Solve the absolute value or indicate that is has no solution.
ex. }$|x+2|-2=10$
\begin{enumerate}
\item Get the absolute value by itself.
\begin{eqnarray*}
\left|x+2\right|-2 & = & 10\\
\left|x+2\right| & = & 12
\end{eqnarray*}
\item Solve for $+$ and $-$.
\begin{eqnarray*}
x+2=-12 & \mbox{ and } & x+2=12\\
x=-14 &  & x=10
\end{eqnarray*}
\item Check! $\left|-14+2\right|-2=10\checkmark$ and $\left|10+2\right|-2=10\checkmark$
\end{enumerate}
\textbf{Solve by factoring. ex. }$75x-25=3x^{3}-x^{2}$
\begin{enumerate}
\item Set the equation equal to $0$ and write in descending order. 
\[
-3x^{3}+x^{2}+75x-25=0
\]
\item Factor by grouping since there are four unlike terms.
\begin{eqnarray*}
-x^{2}(3x-1)+25(3x-1) & = & 0
\end{eqnarray*}
\item Factor and solve for each factor.
\begin{eqnarray*}
(3x-1)(25-x^{2}) & = & 0\\
(3x-1)(5-x)(5+x) & = & 0\\
 & \Rightarrow & x=\nicefrac{1}{3},-5,5
\end{eqnarray*}
\end{enumerate}
\textbf{Find the $x$ and $y$ intercepts from a graph. }
\begin{enumerate}
\item Find all the times the graph intersects the $x-$axis. These are your
$x-$intercepts.
\item Find all the times the graph intersects the $y-$axis. These are your
$y-$intercepts.

\begin{enumerate}
\item [\textbf{ex.}]
\item []\vspace{-.25cm}\renewcommand{\xmax}{10}
\renewcommand{\xmin}{-10}
\renewcommand{\xstep}{0} %must be xmin+n
\renewcommand{\ymax}{10}
\renewcommand{\ymin}{-10}
\renewcommand{\ystep}{0} %must be ymin+n
%%%%%%%
\begin{tikzpicture} 
\begin{axis}[height=5cm, 
	width=5cm, 
	axis lines=middle,
	minor x tick num={4},
	minor y tick num={4},
	grid=both,
	%xtick={0,50,...,250},
	%ytick={0,10,20,...,40},
	%axis line style={-latex}, 
	xmin=\xmin,xmax=\xmax, 
	ymin=\ymin,ymax=\ymax,
	%restrict y to domain=-1:10,
	%no markers, 
	xlabel={$x$},ylabel={$y$}, 
	%every axis x label/.append style={anchor=north}, 
	%every axis y label/.append style={anchor=east}
	]
	
	%\addplot[color=black,solid,mark=*] coordinates {(0,1)} node[right] {\footnotesize{$(0,1)$}};

	\addplot[domain=\xmin:\xmax,thick,samples=100,draw=red,<->]{1/4*(x-3)*(x+4)*(x-1)} node[pos=.525,right] {$f(x)$};
	
\end{axis}
\end{tikzpicture}
\end{enumerate}
\end{enumerate}
\textbf{Find the Domain and the Range.}
\begin{enumerate}
\item Domain$=x's$; Range$=y's$.
\item From a table or list, just find the $x's$ and $y's$. ex. 
\begin{eqnarray*}
\left\{ \left(1,2\right),\left(0,3\right),\left(-1,4\right),\left(-2,5\right)\right\} \\
\mbox{domain=} & \left\{ 1,0,-1,-2\right\} \\
\mbox{range=} & \left\{ 2,3,4,5\right\} 
\end{eqnarray*}
\item From a graph, look at the horizontal values for the domain ($x's$)
and the vertical values for the range ($y's$). If they appear to
go of the graph, then use $\pm\infty$. 

\begin{enumerate}
\item [\textbf{ex.}]
\item []  \vspace{-.25cm}\renewcommand{\xmax}{3.5}
\renewcommand{\xmin}{-3.5}
\renewcommand{\xstep}{0} %must be xmin+n
\renewcommand{\ymax}{2}
\renewcommand{\ymin}{-2}
\renewcommand{\ystep}{0} %must be ymin+n
%%%%%%%
\begin{tikzpicture} 
\begin{axis}[height=5cm, 
	width=5cm, 
	axis lines=middle,
	minor x tick num={4},
	minor y tick num={4},
	grid=both,
	%xtick={0,50,...,250},
	%ytick={0,10,20,...,40},
	%axis line style={-latex}, 
	xmin=\xmin,xmax=\xmax, 
	ymin=\ymin,ymax=\ymax,
	%restrict y to domain=-1:10,
	%no markers, 
	xlabel={$x$},ylabel={$y$}, 
	%every axis x label/.append style={anchor=north}, 
	%every axis y label/.append style={anchor=east}
	]
	
	%\addplot[color=black,solid,mark=*] coordinates {(0,1)} node[right] {\footnotesize{$(0,1)$}};

	\addplot[domain=\xmin:\xmax,thick,samples=200,draw=red]{sin(deg(x)*1.5708)} node[pos=.65,above] {$f(x)$};

\end{axis}
\end{tikzpicture}
\end{enumerate}
\end{enumerate}
\textbf{Determine if the relation is a function. ex. }$\left\{ \left(1,-2\right),\left(2,-1\right),\left(2,0\right),(3,1)\right\} $
\begin{enumerate}
\item Check that each $x$ , in $\left(x,y\right)$, is paired with only
one $y$ value.
\begin{eqnarray*}
1 & \rightarrow & -2\\
2 & \rightarrow & -1\\
2 & \rightarrow & 0\\
3 & \rightarrow & 1
\end{eqnarray*}
\item If any $x$ is paired with two different $y's$ then the relation
is NOT a function.
\begin{eqnarray*}
2 & \rightarrow & -1\\
2 & \rightarrow & 0
\end{eqnarray*}
 So this is NOT a function.
\item Be suspicious if an $x$ appears more than once.
\end{enumerate}
\newpage\textbf{Find the slope passing through the given points.
ex,} $(1,3)$ and $(0,5)$
\begin{enumerate}
\item Use the formula, $m=\frac{y_{2}-y_{1}}{x_{2}-x_{1}}$. Identify $(x_{1},y_{1})$
and $(x_{2},y_{2})$.
\[
x_{1}=1,\mbox{ }y_{1}=3\mbox{ and \ensuremath{x_{2}=0,\mbox{ }y_{2}=5}}
\]
\item Plug the values into the formula.
\[
m=\frac{5-3}{0-1}=\frac{2}{-1}=-2
\]
\end{enumerate}
\textbf{Use the information to write an equation for the line in point-slope
and slope-intercept form. ex. }Slope of $-2$, passing through $(1,3)$.
\begin{enumerate}
\item Find the slope: Use the formula, $y-y_{1}=m(x-x_{1})$. Identify $m$,
$x_{1},$and $y_{1}$. \\
\[
\mbox{The slope is \ensuremath{m=2}, and the point gives \ensuremath{x_{1}=1}and \ensuremath{y_{1}=3}}.
\]
\item Plug the values into the formula.
\[
y-3=-2(x-1)\mbox{ point-slope form}
\]
\item Solve for $y$ to get $y=mx+b$.
\begin{eqnarray*}
y & = & -2x+2+3\\
y & = & -2x+5\mbox{ slope-intercept form}
\end{eqnarray*}
\end{enumerate}
\textbf{Use the information to write an equation for the line in point-slope
and slope-intercept form. ex. }The line perpendicular to $y=-\nicefrac{1}{4}x+2$
and passing through the point $(1,-1)$.
\begin{enumerate}
\item Find the slope: $m=-(-\nicefrac{4}{1})=4$. \\
It is perpendicular so take the negative reciprocal\@. (If it was
parallel we would just take the slope, $m=-\nicefrac{1}{4}$).
\item Plug the values into the formula. \\
$y-(-1)=4(x-1)\Rightarrow y+1=4(x-1)\mbox{ point-slope form.}$
\item Solve for $y$ to get $y=mx+b$.
\begin{eqnarray*}
y & = & 4x-4-1\\
y & = & 4x-5\mbox{ slope-intercept form}
\end{eqnarray*}
\end{enumerate}
\textbf{Evaluate a function at a given value. ex. $f(x)=3x^{2}-x$
}find $f(-2)$ and $f(x+1)$.
\begin{enumerate}
\item Plug the given value where there is an $x$.
\begin{eqnarray*}
f(-2) & = & 3(-2)^{2}-(-2)\\
 & = & 3\cdot4+2=14
\end{eqnarray*}
\begin{eqnarray*}
f(x+1) & = & 3(x+1)^{2}-(x+1)\\
 & = & 3(x^{2}+2x+1)-x-1\\
 & = & 3x^{2}+5x+2
\end{eqnarray*}
\end{enumerate}
\textbf{Determine if a function is even, odd, or neither. ex }$f(x)=3x^{2}-x$,
$g(x)=5x^{3}-x$, and $h(x)=4x^{2}-4$.

Look at the graph to check symmetry or check algebraically. 

$f(x)$ is neither since $3x^{2}$ is even but $-x$ is odd.

$g(x)$ is odd and $h(x)$ is even since each term is odd and even
respectively.
\end{document}
